\section*{Abstract}

\addcontentsline{toc}{section}{Abstract}

\subsection*{Background}
Epilepsy affects approximately 60 million people worldwide. One third of patients remain drug-resistant, and 69\% of deaths from generalized tonic-clonic seizures could be prevented with timely detection \parencite{sveinssonClinicalRiskFactors2020}. EEG-based prediction is effective but unsuitable for ambulatory use due to operational complexity. Wearable devices using autonomic and motion signals offer an alternative, but translating deep learning approaches from controlled settings to real-world deployment remains challenging.

\subsection*{Methods}
Following the PRISMA guidelines and the Webster \& Watson concept matrix framework, this systematic review synthesizes 13 primary studies published between 2013 and 2026. Nine studies address detection and four address forecasting. The studies were mainly evaluated based on algorithmic architecture, biosignal modality, validation methodology, sensitivity, and false alarm rate.

\subsection*{Results}
Convulsive seizure detection achieves high sensitivity (71.6-100\%) primarily through accelerometry. Only two systems achieve a FPR deemed acceptable in ILAE Phase 3 of 0.008-0.028/h: \textcite{spahrDeepLearningbasedDetection2025} (96\% sensitivity, 0.0054/h) and \textcite{fineDetectionKeyAutomated2025} (100\% sensitivity, 0.023/h). Although these studies meet the performance standards of Phase 3, they are not truly Phase 3 because they were only validated in EMU. ECG-based detection of focal seizures fails clinically due to excessive false alarm rates \parencite{reintjesECGBasedDetectionEpileptic2025}.

Seizure forecasting currently faces a performance ceiling. Across methods and modalities, AUC consistently plateaus around 0.74-0.75 \parencite{meiselMachineLearningWristband2020,vielufSeizureMonitoringCombined2025}. Under leave-one-subject-out validation, only 43\% of patients achieve better-than-chance performance. The patient-specific models achieve 100\% success but require weeks to months of individual data, creating a cold-start problem for newly diagnosed patients.

Multi-modal approaches dominate research (69\% of studies) yet the best detection result uses single-modality accelerometry. This suggests biological specificity matters more than sensor count. Wrist-worn devices predominate (62\% of studies). The validation rigor is generally insufficient, only 3 of 13 studies employ prospective designs, and 11 of 13 validate in hospital settings that likely underestimate real-world false alarm rates.

\subsection*{Conclusions}
Wearable seizure detection using DL architectures is close to reaching clinical viability for nocturnal convulsive seizures, but translation to home settings is missing. Forecasting remains immature without established benchmarks or regulatory pathways. Three gaps impede progress regarding forecasting: lack of large-scale home validation, absence of standardized forecasting benchmarks, and no viable solutions yet for approximately 81\% of patients with focal onset seizures \parencite{ayubNaturalHistoryGeneralized2020}. In most studies False alarm rates remain the primary barrier for daytime detection. Future research should prioritize prospective validation in the environment they will be deployed in, address the cold-start problem through adaptive personalization and identify novel biomarkers for focal seizures without motor manifestations.

